\documentclass[12pt,a4paper]{article}

\usepackage[T1]{fontenc}
\usepackage[utf8]{inputenc}
\usepackage[brazil]{babel}
\usepackage{graphicx}
\usepackage{geometry}
\usepackage{listings}
\usepackage{listingsutf8}
\usepackage{xcolor}
\usepackage{setspace}
\usepackage{float}
\usepackage{tikz}
\usepackage{tabularx}

\usetikzlibrary{positioning, shapes.geometric}
\onehalfspacing
\geometry{margin=2.5cm}

\lstset{
    language=SQL,
    basicstyle=\ttfamily\small,
    keywordstyle=\color{blue},
    commentstyle=\color{gray},
    stringstyle=\color{red},
    breaklines=true,
    showstringspaces=false
}

\newcommand{\montarinicio}[3]{
    % =========================
    % CAPA
    % =========================
    \begin{titlepage}
        \centering

        {\large Faculdade Anhanguera}\\[0.5cm]
        {\large Tecnólogo em DevOps}\\[3cm]

        {\Large \textbf{#1}}\\[2cm]

        {\large Gabriel Blank Stift Mousquer}\\[1cm]

        {\large Ivoti/RS}\\
        {\large 2026}

    \end{titlepage}

    % =========================
    % FOLHA DE ROSTO
    % =========================
    \begin{titlepage}
        \centering

        {\large Faculdade Anhanguera}\\[0.5cm]
        {\large Tecnólogo em DevOps}\\[2cm]

        {\Large \textbf{#1}}\\[1.5cm]

        \begin{flushright}
            Trabalho apresentado à disciplina de \\
            \textbf{#2} \\
            Professor: #3 \\
        \end{flushright}

        \vfill

        {\large Gabriel Blank Stift Mousquer}\\
        Matrícula: 2024201126\\

        \vfill

        {\large Ivoti/RS}\\
        {\large 2026}

    \end{titlepage}

    % =========================
    % SEÇÃO: TOC
    % =========================
    \tableofcontents
    \newpage
}


\begin{document}

\montarinicio
    {VETORES, MATRIZES, STRUCTS E PONTEIROS}
    {Algoritmos e Programação Estruturada}
    {Anderson Emídio De Macedo Gonçalves}

\section{Codificação completa}

O código-fonte desenvolvido em linguagem C é apresentado a seguir.

\lstset{language=C}

\lstinputlisting[
    caption={Código em C},
    inputencoding=utf8/latin1,
    basicstyle=\ttfamily\scriptsize
]{./static/main.c}

\newpage

\section{Explicação do algoritmo}

O programa tem como objetivo armazenar 5 valores inteiros em um vetor, calcular a soma desses valores e exibi-los individualmente junto com o resultado final.

Inicialmente, é definida a constante \texttt{qtde\_elementos} com valor 5, que determina o tamanho do vetor. Em seguida, é declarado o vetor \texttt{vetor} com essa quantidade de posições e a variável \texttt{soma}, inicializada com zero para armazenar o total acumulado.

O primeiro laço \texttt{for} é responsável pela leitura dos valores. A cada iteração, o programa solicita ao usuário um número inteiro e o armazena na posição correspondente do vetor. A leitura é feita utilizando \texttt{scanf} com aritmética de ponteiros (\texttt{vetor+i}), que equivale a acessar o endereço da posição \texttt{i} do vetor. Após a leitura, o valor é imediatamente somado à variável \texttt{soma}.

O segundo laço \texttt{for} percorre o vetor novamente, exibindo cada elemento armazenado no formato \texttt{vetor[i] = valor}, garantindo que todos os valores inseridos sejam apresentados individualmente.

Ao final, o programa exibe o valor total armazenado na variável \texttt{soma}, que corresponde à soma de todos os elementos do vetor.

A implementação garante que:
\begin{itemize}
    \item O vetor armazene corretamente 5 valores inteiros;
    \item Os valores sejam inseridos e acessados sem erros;
    \item A soma seja calculada durante a leitura dos dados;
    \item Todos os elementos do vetor sejam exibidos individualmente;
    \item O resultado final da soma seja apresentado corretamente.
\end{itemize}

\newpage

\section{Execução do código}

Foram realizados testes com diferentes conjuntos de valores para validar o funcionamento do programa.

\subsection{Caso 1: 4, 3, 2, 1, 1}

\begin{figure}[H]
    \centering
    \includegraphics[width=0.8\textwidth]{./static/img-1.png}
\end{figure}

\subsection{Caso 1: 40, 30, 20, 10, 10}

\begin{figure}[H]
    \centering
    \includegraphics[width=0.8\textwidth]{./static/img-2.png}
\end{figure}

\subsection{Caso 1: 1, -3, 2, -5, 6}

\begin{figure}[H]
    \centering
    \includegraphics[width=0.8\textwidth]{./static/img-3.png}
\end{figure}

\end{document}
